\section{Pair dynamics and the quantity to estimate}
\label{sec:geometry}

Write \(H_\lambda\) for the Newman deformation in the physical height
normalization of Rodgers and Tao \cite{RodgersTao}. Its full theta source
is specified in Section~\ref{sec:theta}. For a chosen starting slice
\(\lambda_s\), put \(h(\tau,z)=H_{\lambda_s-\tau}(z)\), so \(h_\tau=h_{zz}\).
The target is survival of two specified simple real zeros on
\(0\le\tau\le\Delta\), corresponding to
\([\lambda_s-\Delta,\lambda_s]\) in Newman time. The initial analysis
uses only the local heat equation and real symmetry. The survival
theorem will allow other nonreal zeros.

On a time interval where two real roots \(a<b\) are simple and the
deflated factor is nonzero on their segment, write
\[
h(\tau,z)=(z-a)(z-b)U(\tau,z),\qquad
g=\frac{U_z}{U},\qquad
\delta=b-a,\quad q=\delta^2,\quad c=\frac{a+b}{2}.
\]
Since a simple root \(r\) has velocity \(-h_{zz}(r)/h_z(r)\),
evaluation of the two product derivatives gives
\[
a'=\frac{2}{\delta}-2g(a),\qquad
b'=-\frac{2}{\delta}-2g(b).
\]
Consequently, with \(d=\delta/2\),
\begin{equation}
\label{eq:gap-law}
c'=-g(a)-g(b),\qquad q'=-8-4qB,\qquad
B=\frac{g(b)-g(a)}{\delta}
  =\frac12\int_{-1}^1 g_z(\tau,c+dv)\,\dd v .
\end{equation}
The mutual term is exactly \(-8\) in the squared gap. The remaining
term uses the full deflated factor, including any analytic contribution
to it. These local identities require no infinite zero-motion sum.

If \(C(t)=\int_0^t B(s)\,\dd s\), then integration gives
\begin{equation}
\label{eq:gap-integral}
q(t)=e^{-4C(t)}\left(q(0)-8\int_0^t e^{4C(s)}\,\dd s\right).
\end{equation}
A certified lower bound for \(q(0)\) and an upper envelope for \(C(s)\)
at all preceding times can keep the bracket positive, provided the pair
stays in the region where the estimates hold. An endpoint estimate for
\(C(\Delta)\) leaves the integral undetermined. We use a pointwise upper
force bound below and prove the required confinement separately.

\subsection{A nondegenerate event and its two descriptions}

Suppose \(h\) is jointly analytic near a real double zero at
\((\tau_*,x_*)\), with \(h_{zz}(\tau_*,x_*)\ne0\). In
\(s=\tau-\tau_*\), \(w=z-x_*\), local preparation gives a nonvanishing
factor times \(w^2+A_*(s)w+C_*(s)\), with \(A_*(0)=C_*(0)=0\).
At the event the unit is \(h_{zz}/2\), so the heat equation gives
\(C_*'(0)=2\). The discriminant is therefore \(-8s+O(s^2)\).
The local pair is real and simple for \(s<0\) close to zero and
conjugate for \(s>0\), with upper imaginary part
\(\sqrt{2s}+O(s^{3/2})\). The squared gap extends as the discriminant
through the double zero. Along the real side the factor \(U\) tends to
the analytic quotient by \((z-x_*)^2\). Thus \(B\) has the limit
\begin{equation}
\label{eq:event-curvature}
B_*=\frac{h_{zzzz}}{6h_{zz}}-\frac{h_{zzz}^2}{9h_{zz}^2}
    =\frac{\gamma_*}{6}+\frac{\beta_*^2}{18},\qquad
\beta_*=\frac{h_{zzz}}{h_{zz}},\quad
\gamma_*=\frac{h_{zzzz}}{h_{zz}}-\beta_*^2 ,
\end{equation}
where all derivatives are at the event. Indeed the factor and its
first two derivatives there are \(h_{zz}/2,h_{zzz}/6,h_{zzzz}/12\).

The other description follows the critical point \(x_c(\tau)\) given by
\(h_z(\tau,x_c)=0\) and the implicit function theorem. Its velocity and
critical value \(v_c=h(\tau,x_c(\tau))\) satisfy
\[
x_c'=-h_{zzz}/h_{zz},\qquad v_c'=h_{zz}(\tau,x_c(\tau)).
\]
These equations explain the event-jet alternative. A curvature ratio
at the terminal double root gives a local coefficient, while either
a critical-value comparison or a gap comparison needs bounds over the
preceding interval. For example, on a tracked minimum,
\(v_c(0)+M_{xx}\Delta<0\) and \(h_{zz}(\tau,x_c)\le M_{xx}\) keep
the minimum negative. Positive wall values and strictly signed spatial derivatives
on each side (or strict convexity) can then certify two simple roots.
Any comparison with that method must pay for those geometric data.
The inherited collision and primitive derivations are recorded in the
project sources \cite{OARecord}.
